\chapter{Integrácie}

\section{Integrácie so zariadením}

\subsection{Kamera a galéria}
Aplikácia využíva knižnicu \texttt{react-native-image-picker} na prístup ku kamere a galérii. Pred prvým použitím si aplikácia vyžiada príslušné oprávnenia (\texttt{CAMERA}, \texttt{READ\_MEDIA\_IMAGES} na Android; \texttt{NSCameraUsageDescription}, \texttt{NSPhotoLibraryUsageDescription} na iOS). Fotografie sú odoslané na server ako \texttt{multipart/form-data} bez Base64 kódovania.

\subsection{GPS / Poloha}
Knižnica \texttt{react-native-geolocation-service} zabezpečuje prístup k polohe zariadenia. Worker pri aktívnej práci zdieľa svoju polohu každých 10 sekúnd cez WebSocket. Aplikácia si vyžiada oprávnenia \texttt{ACCESS\_FINE\_LOCATION} (Android) a \texttt{NSLocationWhenInUseUsageDescription} / \texttt{NSLocationAlwaysUsageDescription} (iOS). Na mape (\texttt{react-native-maps} s Google Maps providerom) sú zobrazené live piny workerov a statické piny jobov.

\subsection{Push notifikácie}
Push notifikácie sú implementované pomocou \textbf{Firebase Cloud Messaging} cez knižnicu \texttt{@react-native-firebase/messaging}. FCM token zariadenia je uložený na backende (tabuľka \texttt{users.fcm\_token}) a aktualizovaný pri každom prihlásení. Notifikácia o novom jobe (odoslaná serverom) obsahuje \texttt{data} payload s \texttt{job\_id}, čo umožňuje presmerovanie na Job Detail po kliknutí.

\section{Integrácie s externými službami}

\subsection{PVP4 — Google SSO}
Autentifikácia cez Google je implementovaná pomocou \texttt{@react-native-google-signin/google-signin} a Google OAuth 2.0. Po úspešnom Google Sign-In aplikácia odošle \texttt{id\_token} na backend endpoint \texttt{POST /api/auth/google}, kde je token overený a vydaný JWT.

\subsection{PVP5 — Firebase Analytics + Crashlytics}
\begin{itemize}
  \item \textbf{Firebase Analytics} (\texttt{@react-native-firebase/analytics}) — sleduje kľúčové udalosti: \texttt{job\_created}, \texttt{job\_status\_changed}, \texttt{photo\_uploaded}, \texttt{chat\_message\_sent}. Výstupy z analytiky sú súčasťou prílohy dokumentácie.
  \item \textbf{Firebase Crashlytics} (\texttt{@react-native-firebase/crashlytics}) — automaticky zachytáva pády aplikácie a nekritické chyby. Do dokumentácie je priložený report crashov z testovacej fázy.
\end{itemize}

\subsection{PVP6 — Location Based Service}
Okrem live sledovania polohy je implementovaný \textbf{geofencing}: keď sa Worker vzdiali o viac ako 500\,m od GPS súradníc priradeného jobu, aplikácia zobrazí upozornenie. Geofencing je implementovaný pomocou \texttt{react-native-geofencing}.

\subsection{React Native Maps — Google Maps}
Mapa je renderovaná pomocou \texttt{react-native-maps} s Google Maps providerom. Zobrazuje:
\begin{itemize}
  \item farebné markery polohy každého aktívneho Workera (aktualizované cez WebSocket),
  \item pin miesta každého jobu s farebným rozlíšením podľa stavu (\texttt{pending} = šedá, \texttt{in\_progress} = oranžová, \texttt{done} = zelená).
\end{itemize}